% 6.模型检验.tex —— 六、模型检验
\section{模型检验}
为验证模型与数值格式的有效性，进行以下检验。

\subsection{误差分析}
\textbf{（1）与独立解析解对比。} 构造无限长圆柱、常环境温度对流边界下的解析解（分离变量级数解，Bessel 函数展开，取 60 项）作为基准。取 $T_\infty=50$ $^\circ$C、$N=160$、$\Delta t=0.05$ s，数值解与解析解的最大偏差在 $10^{-5}$ $^\circ$C 量级（$t=300$ s 时 $5.31\times10^{-5}$ $^\circ$C），相对误差小于 $10^{-6}$。该算例同时检验了中心 L'H\^opital 处理与表面 Robin 离散的正确性。

\textbf{（2）空间与时间收敛阶。} 空间方向（固定 $\Delta t=0.02$ s，以 $N=640$ 为参考解）观测阶为 $2.01,2.02,2.07$，温度与浓度场均呈\textbf{二阶收敛}；时间方向（固定 $N=40$）CN 格式观测阶为 $2.00$，呈\textbf{时间二阶}；而 Euler 显式格式时间收敛阶为 $1.00$（一阶）。这正是引入 CN 格式的精度收益。

\textbf{（3）全局守恒回代。} 以守恒权重 $w_0=\Delta r^2/8$、$w_j=r_j\Delta r$、$w_N=R\Delta r/2-\Delta r^2/8$ 作加权和，$\sum_j w_j=\Delta r^2/8$ 精确等于圆柱截面积积分，且所有内部界面通量两两相消，仅余表面项。全局热平衡与质量平衡的相对偏差分别为 $2.93\times10^{-14}$、$1.94\times10^{-14}$，在浮点舍入水平，说明离散守恒为恒等式。

\subsection{灵敏度分析}
\textbf{（1）稳定性对比。} 环境恒为 $T_\infty=50$ $^\circ$C（精确解恒有 $T\le50$ $^\circ$C），任何超界取值均为非物理放大。Crank--Nicolson 格式在 $\Delta t$ 放大至 $300$ s（$300$ 倍）时表面温度仅变化 $1.1$ $^\circ$C，验证\textbf{无条件稳定}；Euler 显式格式在 $Fo=0.405$（$\Delta t=2.4$ s）时有界，而在 $Fo=0.422$（$\Delta t=2.5$ s）时即发散至 $10^6$ 量级，与式 \eqref{eq:cfl} 的临界值 $Fo_{\text{crit}}=0.413$ 精确一致。这定量说明了“由显式升级为 CN”的直接收益：时间步可放大两个数量级而不失稳。

\begin{table}[H]
  \centering
  \caption{显式与隐式格式的稳定性对比（$N=20$，$t_{\text{end}}=1800$ s）}
  \label{tab:stability}
  \resizebox{\textwidth}{!}{%
  \begin{tabular}{lccccc}
    \toprule
    格式 & $\Delta t$ / s & $Fo=\alpha\Delta t/\Delta r^2$ & $T_{\max}$ / $^\circ$C & 物理有界 \\
    \midrule
    \multirow{3}{*}{Euler 显式} & 1.0 & 0.169 & $4.6642\times10^{1}$ & $\checkmark$ \\
     & 2.4 & 0.405 & $4.6645\times10^{1}$ & $\checkmark$ \\
     & 2.5 & 0.422 & $1.1553\times10^{6}$ & $\times$ \\
    \midrule
    \multirow{3}{*}{Crank--Nicolson} & 10 & — & $46.6396$ & $\checkmark$ \\
     & 60 & — & $46.6264$ & $\checkmark$ \\
     & 300 & — & $46.4007$ & $\checkmark$ \\
    \bottomrule
  \end{tabular}}
\end{table}

图 \ref{fig:stability} 直观对比了二者的稳定性行为。

\begin{figure}[H]
  \centering
  \includegraphics[width=0.96\textwidth]{fig8_稳定性对比.pdf}
  \caption{Euler 显式与 Crank--Nicolson 的稳定性对比：（a）显式格式在 $Fo>0.41$ 后非物理发散（对数刻度）；（b）CN 格式任意 $\Delta t$ 保持有界（线性刻度）。本图用于\textbf{数值格式}验证（显式与隐式的稳定性行为）}
  \label{fig:stability}
\end{figure}

\textbf{（2）热物性参数敏感性。} 热扩散系数 $\alpha$ 变化 $\pm20\%$ 使中心温度变化约 $\mp0.6$ $^\circ$C、表面温度仅约 $\pm0.2$ $^\circ$C，且符号相反——$\alpha$ 越小表面热量越难向内传导，表面温度越接近环境温度。故热物性不确定性主要影响内部温度而非表面温度。

\textbf{（3）扩散系数非线性影响。} 将 $D$ 强制取常值 $D(C_0)$ 重算作为基线，常系数近似高估整体失水量 $1.25\%$（表面浓度高估 $2.72\%$）。方向与 $D(C)$ 的负反馈（自减速）机制一致，但幅度小，因 30 min 内浓度变化集中于表层。

\subsection{问题三长时间推进的数值检验}
\label{sec:p3_verify}

问题三的时间跨度是问题二的 $19$ 倍，且终止判据取全节点最大值，故需专门检验。下列检验全部由 \texttt{p3\_verify.py} 自动执行，原始输出见 \texttt{problem3\_verification.txt}（基准网格 $N=20$，参考解为均匀 $\Delta t=1$ s，总耗时约 $7$ min）。

\textbf{（1）时间步无关性（V1）。} 以 $\Delta t = 1$ s 的解为基准，在 $3373$ 个整分钟时刻逐点比对：

\begin{table}[H]
  \centering
  \caption{自适应步长与均匀 $60$ s 步长相对 $\Delta t = 1$ s 参考解的最大偏差（$N=20$，kg/kg）}
  \label{tab:p3_v1}
  \begin{tabular}{lccc}
    \toprule
    口径 & 全程最大偏差 & 第 1 h 内最大偏差 & 优于基准的时刻数 \\
    \midrule
    自适应（60 $\to$ 300 s $+$ 稠密输出） & $3.765\times10^{-6}$ & $3.765\times10^{-6}$ & $3372 / 3373$ \\
    均匀 $60$ s 宏观步                 & $6.267\times10^{-3}$ & $6.267\times10^{-3}$ & $1 / 3373$ \\
    \bottomrule
  \end{tabular}
\end{table}

即自适应格式把时间离散误差压低约 $1.7\times10^{3}$ 倍。机制是：初始瞬变段（表面浓度在 $60$ s 内下降 $0.18$ kg/kg）被自动细分到 $m=32$ 个子步（等效子步长约 $1.9$ s），后段又放回 $300$ s。故该策略同时降低误差与计算量——误差来自瞬变段，算力省在缓变段。

\textbf{（2）稠密输出与端点误差（V2、V4）。} 对步内每个输出节拍点，比较 Hermite 稠密输出与从同一起点以 $\Delta t=1$ s 直接积分的结果，纯插值误差最大 $3.8\times10^{-10}$，比 $4$ 位小数输出所需的 $5\times10^{-5}$ 低 $5$ 个量级。宏观步\textbf{端点}（排除插值）与 $1$ s 参考解的偏差：$h=300$ s 时最大 $6.22\times10^{-7}$（中位数 $1.19\times10^{-8}$，$675$ 个端点中\textbf{无一个}超过 $10^{-3}$）；$h=60$ s 时最大 $3.77\times10^{-6}$（出现在 $t=60$ s）。

\textbf{（3）误差估计与其截断（V3）。} $897$ 个宏观步的误差估计中位数为 $4.2\times10^{-10}$；其中 $264$ 个（$29.4\%$）高于容差 $10^{-8}$，经核查\textbf{全部达到子步上限 $m=32$、全部落在 $t<5.47$ h}。这说明“估计值 $\le$ 容差”并非无条件成立，而是受子步上限截断——这是有意取舍，因为初值到恒温的第一分钟内表面浓度下降 $0.18$ kg/kg，无限细分收益递减。被截断段的\textbf{真实}误差由 V4 独立给出（$3.77\times10^{-6}$），仍比输出要求低一个量级以上。

\textbf{（4）守恒性与求积误差的分离（V5）。} 单步离散恒等式 $\sum_j w_j\Delta C_j=-\Delta t\,[\theta$ 加权表面通量$]$ 的残差为 $9.2\times10^{-9}$，比 Picard 容差 $10^{-10}$ 高约 $2$ 个量级，是 \texttt{picard\_max} $=3$ 截断迭代的余量；因内部界面通量在逐项求和中精确相消，该恒等式由格式结构保证。若改用梯形／Simpson 两种求积格式计算边界通量积分，二者相对偏差分别为 $3.99\times10^{-5}$、$2.55\times10^{-6}$，而“梯形 $-$ Simpson 之差”为 $3.7\times10^{-5}$，与梯形法自身的偏差同量级——说明这 $4\times10^{-5}$ 是 $60$ s 节拍上的 $O(\Delta t^2)$ 求积误差，而非守恒性缺陷。

\textbf{（5）空间与时间不确定性的量级对比（V6、V7）。} 改变时间控制参数（输出节拍 $30/120$ s、步长上限 $k_{\max}=1/5/10$、容差 $10^{-8}/10^{-6}$ 共 $6$ 组）使 $t_{\text{dry}}$ 的偏差全部落在 $\pm0.01$ min 内；而空间网格由 $N=20$ 加密到 $N=160$ 使 $t_{\text{dry}}$ 变化 $0.9$ h。故\textbf{终止时间的不确定性几乎全部来自空间离散}，时间方向不是精度瓶颈。这正是式 \eqref{eq:p3_grid} 边界层加密处理的必要性所在。

\textbf{（6）终止时刻的事后重构校核（V8）。} 冻结形状的准稳态外推（QS-0）在本题末期失效：QS 平衡浓度 $C_{\text{eq}}=0.1423$ 只略低于判据 $0.15$，指数式对其极敏感，外推值由 $t_0=48$ h 时的 $73.71$ h 降到 $54$ h 时的 $56.84$ h，始终高估 $40$ min 以上。计入形状漂移（QS-1，式 \eqref{eq:p3_qs1}）后，自 $36$ h 起的 $5$ 个起点给出\textbf{一致的 $+0.34$ min}。QS-1 基于质量平衡恒等式与形状自相似、采用与二分搜索不同的误差来源，但形状漂移参数由同一数值轨道拟合，故只能作为事后降阶一致性检查；其与二分结果的一致说明终止时间不是某一处实现细节的产物。

\textbf{（7）判据的等价形式（V9）。} 判据 $\max_rC<0.15$ 的等价形式是“各节点达标时刻的最大值”。在 $N=20$、$\Delta r=0.1$ cm 的验证网格上，实测各半径首次达标时刻为：表面 $2.0$ cm 处 $12.58$ h、$1.6$ cm 处 $31.60$ h、$1.2$ cm 处 $43.15$ h、$0.8$ cm 处 $50.55$ h、$0.4$ cm 处 $54.80$ h、中心 $56.18$ h，故\textbf{最迟达标节点恒为中心}。该 $56.18$ h 是粗网格上的中心达标时刻，与网格收敛后的全场烘干时间 $57.17$ h 相差约 $1$ h，二者不可混用。另注：$\arg\max_rC$ 在 $3/3373$ 个时刻取在非中心节点，全部落在 $t\le180$ s 的初始瞬变段，该段最大与次大值之差中位数仅 $1.8\times10^{-15}$（浮点噪声量级），即此处 $C$ 基本是平台值、$\arg\max$ 由舍入决定，不影响判据。

\textbf{（8）分级网格的实现正确性（V10）。} $q=1$ 时新格式的全部算子与问题一、二的均匀实现之差 $\le5.4\times10^{-15}$（双精度舍入水平），网格度量 $(r,r_l,r_r,V,dl,dr)$ 逐位相同（$\max|\Delta|=0$），守恒权重满足 $|\sum_jV_j-R^2/2|=0$ 精确成立。$3\times10^{-5}$ m 表面层内的单元数：均匀网格 $N=640$ 时仍为 $0$（$\Delta r=3.125\times10^{-5}$ m 恰等于层厚），而分级 $q=2$ 的 $N=40$ 为 $1$、$N=60$ 为 $3$。即均匀网格是“跨过”而非“分辨”该薄层，这正是其 $t_{\text{dry}}$ 收敛缓慢（表观阶约 $1.0$）的直接原因。

\textbf{（9）降采样误差（V11）。} 以 $N=1280$ 网格在 $t=55$ h 的剖面为基准：均匀网格的输出点恰为节点，PCHIP 与线性插值给出完全相同的偏差（$N=80$ 时均为 $1.07\times10^{-3}$），即插值不引入额外误差；分级网格节点不对齐，PCHIP 一致优于线性插值（$N=60$ 时 $2.53\times10^{-5}$ 对 $5.48\times10^{-5}$）。交付口径 $N=640$ 的 $\Delta r=3.125\times10^{-5}$ m 恰为 $0.1$ cm 的 $1/32$，节点严格对齐，故 \texttt{result3.xlsx} 的降采样不引入插值误差。

\textbf{（10）一处测量错误的更正（方法论）。} 上述 V1、V4 最初给出“自适应与均匀 $60$ s 的最大偏差均为 $1.03\times10^{-1}$”的结果，据此曾误判该策略“不降低误差、只省算力”。复查发现 $1.03\times10^{-1}$ 是\textbf{测量假象}：当时用 $\operatorname{round}(t/60)$ 按分钟取整匹配参考解，而 $\Delta t=1$ s 的解里 $\operatorname{round}(t/60)=1$ 对应 $t=89$ s（$\operatorname{round}(1.5)=2$），于是“第 $1$ 分钟”与真正的 $60$ s 错位约 $29$ s；在瞬变段 $dC/dt\sim10^{-2}$ /s，该错位恰好伪造出 $1.03\times10^{-1}$。改为精确时刻匹配后，同一测点的真实偏差为 $3.77\times10^{-6}$，V4 的端点误差上界由虚高的 $1.03\times10^{-1}$ 修正为 $6.22\times10^{-7}$。之所以记录此事：该假象的数值大小恰好落在“看起来像合理物理量”的区间，且当时 V1 与 V4 给出同一个可疑数字、互相“印证”。这也是本文对每个关键结论都要求至少两条\textbf{误差来源不同}的证据的原因。

\subsection{问题四移动边界模型的数值检验}
\label{sec:p4_verify}

问题四在移动边界与附录 4 变物性之上，还引入了坐标变换对流项与收缩半径的数值微分，故需专门检验。下列检验全部由 \texttt{p4\_verify.py} 自动执行，原始输出见 \texttt{problem4\_verification.txt}。

\textbf{（1）空间收敛（V1）。} 表 \ref{tab:p4_v1} 显示 $N$ 由 $100$ 增至 $800$ 时 $t_{\text{dry}}$ 的变化：$N=400$ 与 $800$ 结果一致到秒级（均为 $188640$ s $=52.4000$ h），\textbf{空间已收敛}，交付取 $N=800$。这与问题三“空间网格是精度瓶颈”形成对照——归一化固定域网格随边界等比例缩放，不出现收缩末期的边界层跨网格问题。

\begin{table}[H]
  \centering
  \caption{问题四的空间收敛（$t_{\text{dry}}$ vs $N$）}
  \label{tab:p4_v1}
  \begin{tabular}{lccc}
    \toprule
    $N$ & $t_{\text{dry}}$/s & $t_{\text{dry}}$/h & 相邻差/min \\
    \midrule
    100 & 188520 & 52.3667 & --- \\
    200 & 188580 & 52.3833 & $+1.00$ \\
    400 & 188640 & 52.4000 & $+1.00$ \\
    800 & 188640 & 52.4000 & $+0.00$ \\
    \bottomrule
  \end{tabular}
\end{table}

\textbf{（2）时间收敛（V2）。} 表 \ref{tab:p4_v2} 显示 $\Delta t=60$ s 与 $120$ s 均给出 $188640$ s，$\Delta t=30$ s 差 $30$ s（$0.5$ min，在时间离散误差内），\textbf{时间已收敛}。

\begin{table}[H]
  \centering
  \caption{问题四的时间收敛（$t_{\text{dry}}$ vs $\Delta t$，$N=400$）}
  \label{tab:p4_v2}
  \begin{tabular}{ccc}
    \toprule
    $\Delta t$/s & $t_{\text{dry}}$/s & $t_{\text{dry}}$/h \\
    \midrule
    120 & 188640 & 52.4000 \\
    60  & 188640 & 52.4000 \\
    30  & 188610 & 52.3917 \\
    \bottomrule
  \end{tabular}
\end{table}

\textbf{（3）对流项符号的物理检验（V3）。} 表 \ref{tab:p4_v3} 给出对流项取三种符号时的 $t_{\text{dry}}$：$+1$（链式法则严格导出）使 $t_{\text{dry}}$ 由无对流的 $50.83$ h 增大 $1.57$ h 至 $52.40$ h，符合“收缩向内输送水分、减缓干燥”的物理；$-1$ 号使 $t_{\text{dry}}$ 反而减小，非物理。对流项占 $t_{\text{dry}}$ 约 $3\%$，\textbf{不可忽略}。这既验证了符号正误，也说明题面给出的 $-E R'/R$ 是笔误。

\begin{table}[H]
  \centering
  \caption{问题四对流项符号的物理检验（$N=400$）}
  \label{tab:p4_v3}
  \begin{tabular}{lcc}
    \toprule
    对流项符号 & $t_{\text{dry}}$/h & 物理解释 \\
    \midrule
    $+1$（正确） & 52.400 & 收缩向内输送水分 $\to$ 浓缩 $\to$ 减缓干燥 \\
    $0$（无对流） & 50.833 & 仅 $1/R^2$ 缩放 \\
    $-1$（错误） & 48.650 & 收缩向外输送 $\to$ 非物理地加速干燥 \\
    \bottomrule
  \end{tabular}
\end{table}

\textbf{（4）收缩效应（V4）。} 同一套附录 4 物性下，有收缩（附件 2，$R:2\to1.2$ cm）给出 $52.400$ h，无收缩（$R$ 冻结 $2$ cm）给出 $129.083$ h。$R^2$ 缩小 $2.79$ 倍、扩散项 $1/R^2$ 相应放大，$t_{\text{dry}}$ 缩短 $2.46$ 倍，收缩效应显著。

\textbf{（5）判据等价性（V5）。} 全部 $3145$ 个输出时刻中，$\max_r C$ 真实偏离中心（$>10^{-9}$）的时刻为 $0$；另有 $40$ 个时刻的 $\arg\max$ 落在非中心节点，经核查全部发生在 $t\in[60,2520]$ s 的初始段，此时 $C$ 在径向上仍为平台值 $2.55$、中心与内部节点之差仅 $10^{-15}$ 量级（浮点并列噪声）。故\textbf{决定终止时间的是最迟达标节点，即中心}，“各处低于 $0.15$”等价于“中心低于 $0.15$”。

\textbf{（6）退化一致性（V6）。} $R'=0$ 时对流算子 $|\cdot|_{\max}=0$（精确为 $0$），Landau 格式\textbf{逐位退化}为问题一、二的固定域格式，不改变其任何结果。

\textbf{（7）离散质量守恒恒等式（V7）。} 守恒型有限体积 + 表面 Robin 拆分给出离散恒等式 $\frac{\dd M}{\dd t}=-\frac{h_m}{R}(C_N-C_\infty)$（$M=\sum_jV_jC_j$）。前 $2$ h 内数值与解析右端之差 $\max=1.212\times10^{-5}$，即\textbf{质量守恒由格式结构精确保证}（残差来自 Picard 迭代容差，非守恒缺陷）。
